% =====================================================================
%  04  反応工程 — 反応ネットワーク・速度式・触媒失活
%  source_memo: 04_reaction.tex §4.2(反応ネットワークと速度式) /
%               §4.3.2-3(コーキングと活性低下 a(t,T))
%  失活グラフ実データ: pdh_simulator/src/catalyst_model.py _A_MATRIX
%               （要項§3-3 提供データ、温度400-700C × 時間0-30min）
%  方針：反応器モデル（支配方程式・Ergun）は載せない。
%        上=反応式と速度式／下左=触媒失活の要点／下右=a(t,T) グラフ。
%        色は反応式の化学種と、温度で意味を持つグラフの系列のみ。
% =====================================================================
% 活性 a を緑の太字で統一表示するマクロ
\definecolor{ActGreen}{HTML}{1BC24A}
\providecommand{\actA}{}\renewcommand{\actA}{{\color{ActGreen}\scalebox{1.2}{\ensuremath{\boldsymbol{a}}}}}
% 左カラム「触媒失活」見出し＋説明文を丸ごと左右に動かすつまみ。
% 正で右へ・負で左へ（0 で本文左マージン3mmのまま）。見出しと本文は連動して動く。ここの数字だけ変える。
\newlength{\descIndent}\setlength{\descIndent}{-2.8mm}
% 説明文の箱の幅。広げるほどグラフ左の余白を文章が使い、折り返しが減る。
\newlength{\descWidth}\setlength{\descWidth}{46mm}
% グラフ横幅を左から縮める量。大きいほどグラフが小さくなり、文章の余地が増える（右枠はスライド端のまま）。
\newlength{\grShrink}\setlength{\grShrink}{6mm}
\begin{frame}[t]

  \vspace*{0.6em}
  % ===== 上：反応ネットワーク + 速度式（全幅）======================
  {\normalsize
   \renewcommand{\arraystretch}{1.18}%
   \begin{tabular}{@{}l@{\hspace{0.7em}}l@{\hspace{0.6em}}l@{}}
     {\scriptsize 脱水素}
       & \setlength{\fboxsep}{4pt}\setlength{\fboxrule}{1.2pt}%
         \fcolorbox{HeaderBlue}{white}{\boldmath$\mathrm{C_3H_8 \rightleftharpoons {\color{SpRed}C_3H_6} + H_2}$}
       & $\displaystyle r_1 = \actA\,k_1\frac{p_{\mathrm{C_3H_8}} - p_{\mathrm{C_3H_6}}p_{\mathrm{H_2}}/K_{\mathrm{eq}}}{1 + p_{\mathrm{C_3H_6}}/K_B}$ \\
     {\scriptsize クラッキング} & {\boldmath$\mathrm{C_3H_8 \rightarrow {\color{SpBlue}C_2H_4} + CH_4}$}
       & $\displaystyle r_2 = k_2\,p_{\mathrm{C_3H_8}}$ \\
     {\scriptsize 水素化} & {\boldmath$\mathrm{{\color{SpBlue}C_2H_4} + H_2 \rightarrow {\color{SpBlue}C_2H_6}}$}
       & $\displaystyle r_3 = k_3\,p_{\mathrm{C_2H_4}}p_{\mathrm{H_2}}$ \\
     {\scriptsize コーキング} & {\boldmath$\mathrm{{\color{SpRed}C_3H_6} \rightarrow 3\,{\color{SpBlue}CH_{0.5}} + 2.25\,H_2}$}
       & {\scriptsize 物質収支では微量・触媒を失活} \\
   \end{tabular}}

  \vspace{0.12em}
  {\setlength{\fboxsep}{4pt}\setlength{\fboxrule}{0.8pt}%
   \noindent\fcolorbox{HeaderBlue}{white}{%
    \parbox{\dimexpr\linewidth-2\fboxsep-2\fboxrule\relax}{%
      \scriptsize 想定触媒：\textbf{\ce{Cr2O3}-\ce{Al2O3}}、主反応 \(\Delta H^{\circ}_{1}\approx +124\,\mathrm{kJ/mol}\)（吸熱）\par
      \vspace{2pt}
      速度定数は修正アレニウス、基準 $T_0=793\,\mathrm{K}$。$K_B$ はプロピレン吸着で高転化時 $r_1$ を抑制。}}}

  \vspace{0.35em}

  % ===== 下：触媒失活（左・簡潔）と a(t,T) グラフ（右・最大化）====
  %  方針：xラベルを yshift で目盛に密着、図タイトル・出典を負 vspace で詰めて
  %        運転時間の前後・タイトル/出典間の余白を最小化し、グラフ高さを最大化。
  \begin{columns}[T,onlytextwidth]
    \begin{column}{0.26\textwidth}
      \noindent\hspace*{\descIndent}\makebox[0pt][l]{\begin{minipage}[t]{\descWidth}
      \vspace*{1.0em}
      {\bfseries 触媒失活}\par
      \vspace{0.5em}
      {\footnotesize\raggedright\linespread{1.18}\selectfont
        コークが活性点を被覆し失活。\par
        \vspace{0.25em}
        活性 $\actA(t,T)$ は主反応 $r_1$ のみに作用\\
        （副反応には不作用）。\par
        \vspace{0.25em}
        運転温度では数分で大きく低下\par
        \quad $\to$\,\textbf{短周期の再生が不可欠}
      }
      \end{minipage}}
    \end{column}
    \begin{column}{0.74\textwidth}
      \noindent\hspace*{\dimexpr5.90mm+\grShrink\relax}\makebox[0pt][l]{%
      \begin{tikzpicture}
        \begin{axis}[
          font=\scriptsize,
          width=\dimexpr\linewidth-\grShrink\relax, height=5.08cm,
          xmin=0, xmax=30, ymin=0, ymax=1,
          enlarge x limits=false, enlarge y limits=false,
          xtick={0,10,20,30}, ytick={0,0.2,0.4,0.6,0.8,1.0},
          xlabel={運転時間 [min]}, ylabel={活性 $\actA$},
          xlabel near ticks, ylabel near ticks,
          x label style={yshift=2.55mm}, y label style={xshift=1.6mm},
          tick align=inside, xtick pos=both, ytick pos=both,
          legend style={font=\scriptsize, at={(0.985,0.98)}, anchor=north east,
                        draw=none, fill=white, fill opacity=0.7, text opacity=1,
                        inner sep=2pt, row sep=0.4pt},
          legend columns=1, legend cell align=left,
          no markers, line width=0.9pt,
        ]
          \addplot[blue]          coordinates {(0,1)(5,0.94)(10,0.90)(15,0.86)(20,0.84)(25,0.82)(30,0.81)}; \addlegendentry{$400\,^\circ$C}
          \addplot[cyan!70!blue]  coordinates {(0,1)(5,0.87)(10,0.80)(15,0.76)(20,0.72)(25,0.70)(30,0.68)}; \addlegendentry{$450\,^\circ$C}
          \addplot[teal]          coordinates {(0,1)(5,0.80)(10,0.68)(15,0.622)(20,0.57)(25,0.52)(30,0.49)}; \addlegendentry{$500\,^\circ$C}
          \addplot[olive]         coordinates {(0,1)(5,0.67)(10,0.52)(15,0.42)(20,0.36)(25,0.31)(30,0.28)}; \addlegendentry{$550\,^\circ$C}
          \addplot[orange]        coordinates {(0,1)(5,0.52)(10,0.35)(15,0.24)(20,0.18)(25,0.15)(30,0.13)}; \addlegendentry{$600\,^\circ$C}
          \addplot[red!70!black]  coordinates {(0,1)(5,0.38)(10,0.21)(15,0.13)(20,0.09)(25,0.07)(30,0.06)}; \addlegendentry{$650\,^\circ$C}
          \addplot[red]           coordinates {(0,1)(5,0.25)(10,0.12)(15,0.08)(20,0.06)(25,0.05)(30,0.04)}; \addlegendentry{$700\,^\circ$C}
          \node[font=\scriptsize\bfseries, anchor=north, fill=white, fill opacity=0.6,
                text opacity=1, inner sep=1.5pt] at (axis cs:10.5,0.985) {高温ほど急速に失活};
        \end{axis}
      \end{tikzpicture}}\par
      \centering
      \vspace*{-0.92em}
      {\scriptsize 　　　　　触媒活性 $\actA$ の経時低下（温度別）}\par
      \vspace*{-0.66em}
      {\tiny 速度式パラメータ・活性データの出典：第17回プロセスデザイン学生コンテスト要項}
    \end{column}
  \end{columns}

\end{frame}
